\section{Implications for Agentic Verification}
\label{sec:implications}

Each result maps to a requirement in the product specification~\cite{tsaprd}.

\begin{description}
  \item[Overlap is a lead, not a verdict (PRD F4; H1).] Supported. High-overlap pairs are worth investigating, since overlap correlates with kill-set similarity, but a third of them are complementary. The audit reports overlap only as an \emph{inferred} ``investigate'' finding and never as a removal candidate.
  \item[Fine-grained ablation under killed parents (PRD F6a; H3).] Supported, strongly. Most unchecked calls sit inside functions whose whole-body mutants are all killed. The audit must sample call and statement deletions beneath killed functions as well as surviving ones, and must never infer a child result from a parent kill. Deterministic call deletion finds these directly and its population is small enough to run in full; it stays in the MVP.
  \item[Group validation of removals (PRD F7; H4).] Supported, strongly. Every evaluated file lost kills when its leave-one-out-redundant tests were removed together. Leave-one-out results alone cannot authorise combined removal, and greedy covers fitted to one mutant sample need a holdout check before they are shown as more than an upper bound.
  \item[Profile without coverage (PRD F2; H5).] Supported. Coverage, and contexts on top of it, change run time by large and project-specific factors. Timings come from uninstrumented runs; coverage runs are labelled as such.
  \item[Vary order when looking for flakiness (PRD F3; H6).] The only intermittent tests we found were order-dependent: same-order repetition shows one as a consistent failure and never shows the other failing. The audit's deep mode now adds three random-order runs when pytest-randomly is installed, reports tests that fail only in random order as a separate finding, and reports ``not observed in $n$ runs'' with its bound.
  \item[Stable test identifiers (PRD F1).] Not a hypothesis, but a finding of the baseline step: marshmallow parametrizes tests with the current time, so their IDs change on every run and they cannot be tracked across runs. The audit now reports such tests as a separate finding.
  \item[Deletion is not a cheap substitute for classical mutation (PRD F5, F6a; H2).] Not supported, and the plugin changes accordingly. The PRD's feasibility plan treated deterministic deletion as an economical first pass. On these projects it yields fewer survivors per CPU-minute than mutmut's operators, not more. The audit keeps both populations and reports them separately, as the PRD already requires, but no longer presents deletion as the cheaper option. Deletion stays in the MVP for a different reason, given by H3: it is the direct way to find effects that no test checks, and its whole population is small enough to run completely within a deep-mode budget, which makes ``every replacement was tried'' labels possible. Whether its survivors are more often worth acting on is the open question the pending review addresses.

  \item[What the benchmark changed.] The benchmark found the mutant-generation defect described in Section~\ref{sec:benchmark}, which is fixed, and showed that coverage summaries without line numbers sent one agent to reconcile counts by hand; the report now lists uncovered lines. It also showed that a 12-minute deep audit is too short for the larger projects on a busy machine. The skill should state the expected time from the suite size before choosing a budget.
  \item[Verify changes as they are made (PRD V1; H7).] Real commits in mature projects leave mutants of their changed lines uncaught, most often in the larger changes, so a check on each change is useful even where the suite is strong overall. A change with no applicable mutant is reported as incomplete, never as verified.
  \item[Check new tests against the old code without touching the user's tree (PRD V2; benchmark V1).] Most agents, with or without the skill, stashed the user's change to see their tests fail on the old code. The one skill agent that obeyed the rule against editing production code skipped the check and missed the real change. \texttt{base-check} gives that check in a disposable copy, and the skill now requires it.
  \item[Removal sets need held-out validation (PRD V3; H8).] Jointly chosen sets kept every sampled kill but lost \HEightLostShare{} of held-out kills in every project. The audit shows them only as an upper bound, and the skill must run \texttt{prune-check} before suggesting any removal.
\end{description}
